\section{Introduction}\label{sec:intro}

\paragraph{The question.}
A physical system sometimes hides a smaller logical system inside it: the
logical observables are encoded in the physical ones, and a recovery operation
returns the system to the encoded subspace. Mathematically, such a recovery is a
quantum channel $T$ that is idempotent, $T^2=T$. Kitaev~\cite{Kitaev2024} studied
the realistic case where the recovery is only \emph{approximately} idempotent,
and asked what algebraic structure the encoded observables carry. He worked with
the dual map $\Phi$ of the channel, a unital completely positive map on the
algebra $B(\HH)$ of operators on a finite-dimensional Hilbert space, and measured
non-idempotence by $\norm{\Phi^2-\Phi}_{\mathrm{cb}}\le\eta$. His answer: the
near-fixed points of $\Phi$ form an \emph{approximate $C^*$-algebra}, and every
finite-dimensional approximate $C^*$-algebra is, with universal (dimension-free)
constants, close to a genuine $C^*$-algebra.

\paragraph{This report.}
We ask the same question for maps that are only \emph{positive}, not completely
positive. This is the natural setting whenever one cares about distinguishability
of states rather than full quantum operations, and it is the setting of
van~Luijk and Wilming~\cite{vanLuijkWilming2026}, who show that for positive
trace-preserving maps the relevant invariant is a \emph{Jordan} algebra. The
present report does three things, in order of completeness:
\begin{enumerate}[leftmargin=*]
  \item It records the consensus \emph{definitions} (positive map, Jordan
    product, JB-algebra, order-unit space, approximate Jordan algebra). Standard
    background definitions are matched against local source copies; project
    definitions such as the $\eps$-JB notion are marked \status{(consensus)} and
    tied to internal proof files in \texttt{PROVENANCE.md}.
  \item It proves the \emph{bridge theorem} (\Cref{sec:bridge}): from a unital
    positive almost-idempotent $\Phi$ one obtains, with no use of dilations or
    complete positivity, a space that satisfies the Jordan-algebra axioms up to
    $O(\sqrt\eta)$. This is the positive-map analogue of the first half of
    Kitaev's construction, and it is complete.
  \item It states the \emph{structure-theory programme} (\Cref{sec:programme})
    --- the analogue of the second half of Kitaev's theorem --- and reports
    precisely what is established and what is still open
    (\Cref{sec:discussion}).
\end{enumerate}

\paragraph{Why Jordan, and not associative.}
The reason positivity forces a Jordan, rather than associative, structure can be
stated in one line. Kitaev's verification of the $C^*$-norm axiom rests on the
Schwarz inequality $\Phi(x)^*\Phi(x)\le\Phi(x^*x)$, which holds for $2$-positive
(in particular completely positive) maps. A merely positive unital map need not
satisfy it; what it does satisfy is Kadison's inequality
$\Phi(x)^2\le\Phi(x^2)$ for self-adjoint $x$~\cite{Kadison1952}. The latter
involves only self-adjoint elements and squares --- exactly the data on which a
Jordan algebra is built. Replacing the associative product $xy$ by the symmetric
product $x\jp y=\tfrac12(xy+yx)$ throughout is therefore not a convenience but a
necessity, and the entire theory reorganises around it.

\paragraph{Conventions and status flags.}
All Hilbert spaces are finite-dimensional and complex. We write $B(\HH)$ for the
operators on $\HH$ and $\Bsa$ for the self-adjoint ones. Statements carry a
status: \status{(proved)} for results with a complete proof in this report;
\status{(cited)} for results transcribed verbatim from a named source;
\status{(consensus)} for definitions and conventions adopted by agreement
between the authors (a definition is never ``proved''); and \status{(open)} for
targets of the programme not yet established. Constants denoted $C$ are universal and never depend on
$\dim\HH$ unless explicitly stated.
